% Section 4.2, from lectures/L04/README.md: the objectives, the evaluation questions, and the next
% lecture.

\section{Review}
\label{c4:sec:review}

\needspace{6\baselineskip}
\subsection*{What you should be able to do now}
After this chapter, you should be able to:
\begin{itemize}
  \item \textbf{Build a synchronous FIFO} with a look-then-advance read, and explain why that split
    is what lets \chapref{5} keep \code{RX_DATA} a pure read.
  \item \textbf{Say what \code{full} and \code{empty} guard}, and what a write to a full FIFO or a
    read from an empty one must do instead of wrapping the pointers past each other.
  \item \textbf{Instantiate the receive path into \code{uart_top}} in the right order, with the
    synchronizer on the pin and the receiver behind it, and explain why binding \code{rx} straight
    to \code{uart_rx} would analyze cleanly and still be wrong.
  \item \textbf{Explain what an overrun is}, why the receiver cannot detect one on its own, and
    which block has the one piece of information that makes it detectable.
\end{itemize}

\needspace{6\baselineskip}
\subsection*{Questions to test yourself}
You should be able to explain:
\begin{enumerate}
  \item \code{rdata} shows the front entry with no \code{rd} at all. What would break in the
    \code{RX_DATA} and \code{RX_POP} split of \chapref{5} if the FIFO instead popped on every read?
  \item A write arrives while the FIFO is full. What must happen to \code{wdata}, to \code{head} and
    to \code{count}, and what tells the caller it happened?
  \item What is an overrun, why can the receiver not detect one on its own, and which block would
    have to do it instead?
\end{enumerate}

\needspace{6\baselineskip}
\subsection*{Looking ahead}
In \chapref{5} ports become registers: poll-able \code{STATUS} bits, sticky error flags, a
write-triggered TX push, the read-then-pop RX path over two of the FIFOs you just built, and
completing \code{uart_top} behind the provided SPI transport, so that the system testbench finally
runs.
